\scanpage{87}


\uline{Corollaire} C'est aussi la condition néc. et suff. pour que, pour tout Module $M$, on ait

\[ \underline{\operatorname{Hom}}_{\underline O}(E_1*E_2,M)\xrightarrow{\sim}\underline{\operatorname{Hom}}(E_1\otimes E_2,M)^{\sigma_1\otimes\sigma_2} \]

Posant $E_1^\vee=\unsure{\mathcal A_1|A_1}$, $E_2^\vee=\mathcal A_2$, $(E_1*E_2)^\vee\simeq\mathcal A$, cela implique que

\[ \mathcal A\longrightarrow\mathcal A_1\otimes\mathcal A_2 \]

induit un isomorphisme

\[ \mathcal A\xrightarrow{\sim}(\mathcal A_1\otimes\mathcal A_2)^{\sigma_1^\vee\otimes\sigma_2^\vee}, \]

et inversement si on veut que, ce \deleted{soit} reste vrai après tout changement de base.

\begin{center}$\backslash$\end{center}

Complément sur les $\unsure{\gamma|p|\chi}$ trivialisations, au dessus de $X_0=V(\gamma)$ i.e sur un schéma $X$ sur lequel $\gamma1_X=0$.

La donnée d'une $\gamma$-trivialisation pour un \inserted{$L$-}torseur $P$, comme $\gamma P=0.P\simeq L$, équivaut simplement à une section $T$ de $L$,

et $\pi:P\longrightarrow L$ est l'application constante $x\longmapsto\underset{\substack{\Vert\\0}}{\gamma x}-T=-T$.

\[ \pi(x)=-T\overset{\rm def}{=}b \]

\unsure{En tous cas|Un bon cas|Vu ! Dans le cas}

\[ E_1\otimes E_2\mathbin{\unsure{-|\to}}E_1*E_2 \]
% archive sheet number: 86
